\begin{textsample}{POS Dim 1 – 2023\_c5 – Score 22.00 – 2023\_c5/enade\_2023\_eng\_civ\_4a1049be40.txt}  \label{ex:f1_pos_003}
115 Relatório Síntese de Área – ENADE/2023 CAPÍTULO 4 DISTRIBUIÇÃO DE CONCEITOS A \textbf{finalidade} deste capítulo é apresentar um \textbf{panorama} nacional sobre a distribuição do Conceito Enade 2023.
Para traçar esse \textbf{panorama}, recorre -se a tabelas e gráficos, que são acompanhados de comentários e análises que contribuem para a contextualização dos dados neles expostos.
Para \textbf{visualizar} o \textbf{panorama} do Enade 2023, em nível nacional, na primeira seção deste capítulo, os dados sobre o número de cursos participantes são organizados por grande região.
Na \textbf{segunda}, o mapeamento do conceito do exame é feito por categoria administrativa das IES e modalidade de oferta dos cursos, de maneira \textbf{cruzada} às grandes regiões.
Na última seção, a lógica de apresentação dos dados é feita por organização acadêmica e, também, por grande região de funcionamento dos cursos, no país como um todo. 4.1 \textbf{PANORAMA} NACIONAL DA DISTRIBUIÇÃO DOS CONCEITOS A lógica adotada na apresentação dos dados, nesta seção e nas duas seguintes, pressupõe o registro de duas informações de natureza \textbf{reiterativa}.
O primeiro diz respeito aos valores possíveis do Enade, que variam em níveis de 1 a 5, sendo que, à medida que esses valores aumentam, o curso é melhor avaliado, e ao fato de que aqueles cursos que tiveram somente um ou nenhum concluinte participante não obtêm o Conceito Enade, ficando sem conceito ( SC ), como ressaltado na apresentação deste relatório.
O segundo registro é que, caso haja diferença entre os cursos \textbf{tabulados} neste capítulo e no capítulo 1, ela corresponderá aos cursos que não foram avaliados, ou seja, àqueles, em princípio, sem estudantes concluintes que tivessem sido inscritos para a prova.
Na Tabela 4.1, são apresentadas a quantidade e a distribuição de cursos de Engenharia Civil participantes do Enade 2023, por faixa de conceito e grande região.
A partir dos dados da Tabela 4.1, nota -se que, dos 917 cursos participantes, 375 ( 40,9 \% ) classificaram -se com conceito 2, o valor modal.
Esse foi também o conceito modal na maioria das regiões : Centro-Oeste ( 41,9 \% ), Nordeste ( 39,6 \% ), Norte ( 48,5 \% ) e Sudeste ( 44,4 \% ).
O conceito 3 foi o segundo mais frequente em nível nacional ( 33,3 \%, correspondendo a 305 cursos ), tendo sido o segundo conceito mais frequente na maioria das regiões, exceto no Sul.
O conceito 4 foi o terceiro ( 12,3 \%, correspondendo a 113 cursos ).
Houve, ainda, 52 ( 5,7 \% ) cursos que receberam o conceito 5 e 31 ( 3,4 \% ) cursos que receberam o conceito 1.
Dos 917 cursos de Engenharia Civil, 41 ( 4,5 \% ) ficaram sem conceito ( SC ).
A região Norte participou com 68 cursos, o que corresponde a 7,4 \% do total nacional.
Desses, 33 cursos receberam o conceito 2, o que \textbf{equivale} a 48,5 \% do total regional, o conceito modal na região, conforme mencionado.
A 23 cursos ( 33,8 \% ) atribuiu -se o conceito 3 ; a oito cursos ( 11,8 \% ), o conceito 4 ; e a dois cursos ( 2,9 \% ) o conceito 1.
Ainda nessa região, nenhum curso ( 0,0 \% ) obteve conceito 5 e dois cursos ( 2,9 \% ) ficaram sem conceito ( SC ).
A região Nordeste participou com 182 cursos, correspondentes a 19,8 \% do total nacional.
Nessa região, 72 cursos ( 39,6 \% em termos regionais ) obtiveram conceito 2, o conceito modal para a região, como já comentado.
Os conceitos 3 e 4 foram atribuídos, respectivamente, a 48 e a 30 cursos ( 26,4 \% e 16,5 \%, respectivamente ).
Ainda nessa região, 17 cursos ( 9,3 \% ) foram avaliados com o conceito 5 e cinco cursos ( 2,7 \% ) foram avaliados com o conceito 1.
Outros dez cursos ( 5,5 \% ) ficaram sem conceito ( SC ).
Para os 394 cursos participantes da região Sudeste ( 43,0 \% do total nacional ), o conceito 2 foi o conceito modal, como já anteriormente assinalado, tendo sido obtido por 175 cursos ( 44,4 \% ).
O conceito 3 foi atribuído a 124 cursos ( 31,5 \% ) e o conceito 4, a 37 cursos ( 9,4 \% ).
Receberam os conceitos 5 e 1, respectivamente, 23 cursos ( 5,8 \% ) e 18 cursos ( 4,6 \% ).
Nessa região, 17 cursos ( 4,3 \% ) ficaram sem conceito ( SC ).
Os 180 cursos da região Sul corresponderam a 19,6 \% do total nacional.
Como já \textbf{indicado}, houve predominância do conceito 3 ( conceito modal ), atribuído a 42,2 \% dos cursos da região, ou seja, atribuído a 76 dos 180 cursos participantes na região Sul.
O conceito 2 foi atribuído a 56 cursos ( 31,1 \% ) ; o conceito 4, 27 cursos ( 15,0 \% ) ; e o conceito 5, a oito cursos ( 4,4 \% ).
Na região Sul, quatro cursos ( 2,2 \% ) obtiveram conceito 1 e nove cursos ( 5,0 \% ) ficaram sem conceito ( SC ).
Já dos 93 cursos participantes na região Centro-Oeste ( 10,1 \% do total nacional ), 39 ( 41,9 \% em termos regionais ) receberam conceito 2, o conceito modal, como já destacado. 34 cursos ( 36,6 \% ) obtiveram o conceito 3, 11 cursos ( 11,8 \% ) obtiveram o conceito 4 e quatro cursos ( 4,3 \% ), o conceito 5.
Nessa região, dois cursos ( 2,2 \% ) obtiveram o conceito 1 e três cursos ( 3,2 \% ) ficaram sem conceito ( SC ).
Em um gráfico de distribuição \textbf{cumulativa}, em cada valor de \textbf{abscissa} ( eixo x ), o valor apresentado no eixo das \textbf{ordenadas} ( eixo y ) é o valor acumulado da \textbf{variável} ( Conceito Enade ).
Em linhas gerais, em um gráfico desta mesma \textbf{tipologia}, \textbf{poligonais} mais à esquerda representam, no caso dos dados apresentados nesta seção, as grandes regiões com a pior distribuição, e \textbf{poligonais} mais à direita, as grandes regiões com os melhores conceitos.
A partir destas diretrizes, no Gráfico 4.1, é apresentada a distribuição \textbf{cumulativa} do Conceito Enade, segundo a grande região.
No Gráfico 4.1, observando a linha vertical \textbf{pontilhada}, que ajuda a \textbf{visualizar} o percentual de conceitos acumulados até o valor 2, ou seja, conceitos 1 e 2, pode -se notar, por exemplo, que a região Norte ( linha azul ) apresenta 53,0 \% dos cursos com conceito 2 ou abaixo deste valor.
A região Sul ( linha \textbf{roxa} ), por outro lado, apresenta 35,1 \% dos cursos na mesma situação : com conceito 2 ou abaixo. -Destaca -se que, a região Norte ( linha azul ), o conceito 5 não foi alcançado por nenhum curso, logo os cursos têm conceito menor ou igual a 4, e a respectiva \textbf{poligonal} \textbf{cumulativa} alcança 100 \% no conceito 4.
A região Norte ( linha azul ) apresenta os cursos com conceitos mais baixos, ao passo que a região Sul ( linha \textbf{roxa} ) apresenta a distribuição com valores maiores.
Os cursos em IES situadas na região Centro-Oeste ( linha laranja ) apresentam uma situação intermediária. 4.2 CONCEITOS POR CATEGORIA ADMINISTRATIVA, MODALIDADE DE OFERTA DOS CURSOS E GRANDE REGIÃO Na Tabela 4.2, é apresentada a distribuição dos cursos de Engenharia Civil participantes do Enade 2023, por categoria administrativa e modalidade de oferta, de acordo com os conceitos por eles alcançados, segundo as grandes regiões brasileiras.
De acordo com as informações da Tabela 4.2, dos 917 cursos participantes, 175 ( 19,1 \% ) eram \textbf{ministrados} em instituições públicas e 742 ( 80,9 \% ), em instituições privadas.
Quanto à modalidade de oferta, 863 ( 94,1 \% ) cursos eram presenciais e 54 ( 5,9 \% ), a distância.
De acordo com as informações da mesma tabela, em termos nacionais, entre instituições públicas e privadas, observa -se que, dos 52 cursos avaliados com conceito 5, 49 foram oferecidos em IES públicas e 3, em IES privadas.
Dos 742 cursos participantes de IES privadas, o conceito 2 foi o valor modal, atribuído a 368 ( 49,6 \% ) cursos.
Em relação aos demais cursos de instituições privadas, 261 cursos ( 35,2 \% ) foram avaliados com o conceito 3, 39 cursos ( 5,3 \% ) foram avaliados com o conceito 4, 31 cursos ( 4,2 \% ) foram avaliados com o conceito 1 e, três cursos ( 0,4 \% ), com o conceito 5.
Nessa categoria, 40 cursos ( 5,4 \% ) ficaram sem conceito ( SC ).
Na rede pública, o conceito modal foi o 4, atribuído a 74 cursos, correspondentes a 42,3 \% dos 175 cursos da categoria.
Dos demais cursos participantes, 49 cursos ( 28,0 \% ) obtiveram o conceito 5, 44 cursos ( 25,1 \% ) receberam o conceito 3 e sete cursos ( 4,0 \% ), o conceito 2.
O conceito 1 não foi atribuído a nenhum curso ( 0,0 \% ).
Nessa categoria administrativa, um curso ( 0,6 \% ) ficou sem conceito ( SC ).
Dos 863 cursos participantes da modalidade presencial, o conceito 2 foi o valor modal, atribuído a 342 ( 39,6 \% ) cursos.
Dos demais cursos dessa modalidade de oferta, 289 cursos ( 33,5 \% ) receberam o conceito 3 ; 112 cursos ( 13,0 \% ) receberam o conceito 4 ; 52 cursos ( 6,0 \% ), o conceito 5 ; e 29 cursos ( 3,4 \% ), o conceito 1.
Nessa modalidade de oferta, 39 cursos ( 4,5 \% ) ficaram sem conceito ( SC ).
Dos 54 cursos participantes da modalidade a distância, o conceito 2 foi o valor modal, atribuído a 33 ( 61,1 \% ) cursos.
Dos demais cursos \textbf{ofertados} nessa modalidade, 16 cursos ( 29,6 \% ) receberam conceito 3, dois cursos ( 3,7 \% ) receberam o conceito 1 ; um curso ( 1,9 \% ) recebeu o conceito 4 ; e nenhum curso ( 0,0 \% ) recebeu conceito 5.
Nessa modalidade de oferta, dois cursos ( 3,7 \% ) ficaram sem conceito ( SC ).
Na análise por região, observa -se, que, na região Norte, as instituições públicas participaram com 18 cursos ( 26,5 \% total regional ), dos quais nenhum ficou sem conceito ( SC ).
O conceito modal foi 3, atribuído a dez cursos, correspondendo a 55,6 \% dessa categoria na região.
O conceito 4 teve uma concentração muito próxima do modal, com sete cursos ( 38,9 \% ), e o conceito 2 teve um curso ( 5,6 \% ).
Dos demais cursos, nenhum recebeu os conceitos 5 e 1.
As instituições privadas participaram com 50 cursos na região Norte ( 73,5 \% do total regional ), e dois cursos ficaram sem conceito ( SC ).
O conceito modal nas instituições privadas na região foi o 2, atribuído a 32 cursos, o que corresponde a 64,0 \% dessa categoria na região.
Os conceitos 3 e 1 foram atribuídos a treze e a dois cursos, respectivamente.
O conceito 4 foi atribuído a um curso e o conceito 5 não foi atribuído a nenhum curso ( 0,0 \% ).
Dos cursos oferecidos na região Norte, 63 eram presenciais e um ficou sem conceito ( SC ).
Dos 63 cursos oferecidos, 30 obtiveram o conceito 2 ( modal ) ; 22, o conceito 3 ; oito, o conceito 4 ; dois, o conceito 1 ; e nenhum, o conceito 5.
Cinco cursos eram na modalidade a distância, tendo três recebido o conceito modal 2 ; um, o conceito 3 ; e nenhum, os conceitos 5, 4 e 1.
Um curso ficou sem conceito ( SC ).
Na região Nordeste, a rede privada concentra 130 dos 182 cursos participantes, o equivalente a 71,4 \% do total da região.
Entre os cursos oferecidos por IES privadas no Nordeste, 71 cursos foram avaliados com conceito 2, o conceito modal. 39 cursos obtiveram o conceito 3 ; cinco cursos, o conceito 4 e o 1 ; e um curso obteve o conceito 5.
Nessa categoria administrativa e grande região, nove cursos ficaram sem conceito ( SC ).
As instituições públicas da região Nordeste participaram com 52 cursos ( 28,6 \% do total da região ).
Desses, 25 ( 48,1 \% ) obtiveram conceito 4, o valor modal, seguidos por 16 cursos, que receberam o conceito 5.
O conceito 3 foi atribuído a nove cursos ; o conceito 2, a um curso ; e o conceito 1, a nenhum curso.
Um curso ficou sem conceito ( SC ).
Dos cursos oferecidos na região Nordeste, 176 eram presenciais : dez ficaram sem conceito ( SC ), quatro obtiveram o conceito 1 ; 67, o conceito 2, o valor modal ; 48, o conceito 3 ; 30, o conceito 4 ; e 17, o conceito 5.
Seis cursos eram na modalidade a distância, tendo cinco cursos recebido o conceito modal 2 ; um, o conceito 1 ; nenhum, os conceitos 3, 4 e 5.
Nenhum curso ficou sem conceito ( SC ).
Na região Sudeste, a proporção de cursos da rede privada, 87,6 \%, correspondeu a 345 dos 394 cursos participantes.
Nessa combinação de categoria administrativa e grande região, o conceito modal foi o 2, atribuído a 172 cursos ; e 17 cursos ficaram sem conceito ( SC ).
Os demais foram avaliados com os seguintes conceitos : 3 ( 117 cursos ), 4 ( 19 cursos ), 1 ( 18 cursos ) e 5 ( dois cursos ).
Dos 49 cursos oferecidos por instituições públicas na região Sudeste, 21 obtiveram o conceito 5 ; 18, o conceito 4 ; sete, o conceito 3 ; três, o conceito 2 ; e nenhum obteve o conceito 1.
Nenhum curso ficou sem conceito ( SC ).
Ainda na região Sudeste, grande parcela dos cursos ( 94,7 \% ) foi oferecida na modalidade presencial : 373 dos 394 cursos.
Para esses, a moda foi o conceito 2, com 165 cursos.
Os demais cursos presenciais da região foram alocados nos conceitos 3 ( 114 cursos ), 4 ( 36 cursos ), 5 ( 23 cursos ) e 1 ( 18 cursos ). 17 cursos presenciais ficaram sem conceito ( SC ).
Na região Sudeste, 21 cursos foram oferecidos, na modalidade a distância, dez com conceito 2 ; dez com conceito 3 ; um com conceito 4 ; e nenhum com conceito 1 nem com 5.
Nenhum curso ficou sem conceito ( SC ).
As instituições privadas concentraram 149 dos 180 cursos participantes da região Sul, 82,8 \% do total regional.
Desses, 70 cursos obtiveram conceito 3, o conceito modal.
Nessa combinação de categoria administrativa e grande região, 56 receberam o conceito 2 ; dez receberam o conceito 4 ; quatro, o conceito 1 ; e nenhum obteve o conceito 5.
Nove cursos ficaram sem conceito ( SC ).
As instituições públicas na região Sul participaram com 31 cursos ( 17,2 \% dos cursos da região ), aos quais foram atribuídos os conceitos 1 e 2 ( nenhum curso ), 3 ( seis cursos ), 4 ( 17 cursos ) e 5 ( oito cursos ).
Nenhum curso da região Sul oferecido por IES públicas ficou sem conceito ( SC ).
Na região Sul, 161 dos cursos foram oferecidos na modalidade presencial.
Para esses, a moda foi o conceito 3, com 72 cursos.
Os demais cursos presenciais da região foram alocados nos conceitos 2 ( 42 cursos ), 4 ( 27 cursos ), 5 ( oito cursos ) e 1 ( três cursos ).
Ficaram sem conceito ( SC ) nove cursos.
Na modalidade a distância, foram oferecidos 19 cursos, sendo um curso com conceito 1, 14 cursos com conceito 2 ( modal ), quatro cursos com conceito 3, nenhum curso com conceito 4 nem com conceito 5.
Nenhum curso ficou sem conceito ( SC ) nessa região.
Na região Centro-Oeste, 68 dos 93 cursos participantes eram de IES privadas ( 73,1 \% em termos regionais ).
Desses, 37 cursos receberam conceito 2, o valor modal.
Dos demais cursos, 22 receberam o conceito 3, quatro receberam o conceito 4, dois receberam o conceito 1 e nenhum recebeu o conceito 5.
Três ficaram sem conceito ( SC ).
Quanto aos 25 cursos oferecidos por instituições públicas na região Centro-Oeste ( 26,9 \% do total regional ), nenhum ficou com conceito 1 ; dois, com conceito 2 ; 12, com conceito 3 ; sete, com conceito 4 ; e quatro, com conceito 5.
Nenhum ficou sem conceito ( SC ) nessa região.
Também na região Centro-Oeste, a maioria dos cursos foi \textbf{ofertada} na modalidade presencial : 90 ( 96,8 \% ).
Para esses, a moda foi o conceito 2, com 38 cursos ( 42,2 \% ).
Os demais cursos presenciais da região foram alocados nos conceitos 3 ( 33 cursos, 36,7 \% ), 4 ( 11 cursos, 12,2 \% ), 5 ( quatro cursos, 4,4 \% ) e 1 ( dois cursos, 2,2 \% ).
Dois cursos presenciais da região ficaram sem conceito ( SC ).
Na modalidade a distância, foram oferecidos três cursos, sendo nenhum curso com conceito 1 ; um curso com conceito 2 ; um, com conceito 3 ; e nenhum curso com conceito 4 nem com conceito 5.
Um curso ficou sem conceito ( SC ) nessa região.
A título de síntese, no Gráfico 4.2, apresenta -se a distribuição do Conceito Enade, segundo a categoria administrativa da IES.
Nesse gráfico, como dito anteriormente, \textbf{poligonais} mais à esquerda representam Conceito Enade mais baixos, e \textbf{poligonais} mais à direita, Conceito Enade mais altos.
É importante observar que, para cada \textbf{poligonal} presente nesse gráfico, a porcentagem constante no eixo y refere -se ao universo de cursos oferecidos por instituições públicas ( \textbf{poligonal} azul ) ou privadas ( \textbf{poligonal} verde ), respectivamente.
Os cursos em IES públicas ( \textbf{poligonal} azul ) apresentam uma \textbf{poligonal} mais à direita do que os em IES privadas ( \textbf{poligonal} verde ) e, consequentemente, uma distribuição de Conceito Enade com valores maiores. 4.3 CONCEITOS POR ORGANIZAÇÃO ACADÊMICA E GRANDE REGIÃO Na Tabela 4.3, encontra -se a distribuição dos conceitos atribuídos aos cursos participantes do Enade 2023, na área de Engenharia Civil, por organização acadêmica, segundo as grandes regiões brasileiras.
Como mostram os dados da Tabela 4.3, dos 917 cursos de Engenharia Civil participantes, 300 eram oferecidos em universidades ; 314, em centros universitários ; 256, em faculdades ; e 47, por CEFET/IF.
Essa distribuição corresponde a, respectivamente, 32,7 \%, 34,2 \%, 27,9 \% e 5,1 \% dos cursos.
De acordo com os dados apresentados na mesma tabela, de todos os 375 cursos avaliados com o conceito 2, 77 eram vinculados a universidades.
Os cursos vinculados a esse tipo de organização acadêmica tiveram o conceito 3 como conceito modal, atribuído a 115 cursos ( 38,3 \% ).
Os demais cursos \textbf{ofertados} por universidades receberam os conceitos 4 ( 60 cursos ), 5 ( 44 cursos ) e 1 ( três cursos ).
Um curso ficou sem conceito ( SC ).
Dos cursos vinculados a centros universitários, o conceito modal também foi o 2, atribuído a 145 cursos ( 46,2 \% ). 22 cursos vinculados a esse tipo de organização acadêmica ficaram sem conceito ( SC ), e os outros cursos receberam os conceitos 3 ( 117 cursos ), 4 ( 17 cursos ), 1 ( 11 cursos ) e 5 ( dois cursos ).
Dos 256 cursos mantidos por faculdades, 152 ( 59,4 \% ) receberam conceito 2, o conceito modal.
O segundo conceito mais frequentemente atribuído aos cursos oferecidos por essa categoria de IES foi o conceito 3, recebido por 58 cursos.
Na sequência, 17 cursos obtiveram conceito 1 ; oito cursos, o conceito 4 ; e três cursos, o conceito 5.
Dos cursos oferecidos em faculdades, 18 ficaram sem conceito ( SC ).
Dos cursos oferecidos por CEFET/IF, o conceito modal foi o 4, atribuído a 28 cursos ( 59,6 \% ).
Nenhum curso vinculado a esse tipo de organização acadêmica ficou sem conceito ( SC ), e os outros cursos receberam os conceitos 3 ( 15 cursos ), 5 ( três cursos ), 2 ( um curso ) e 1 ( nenhum curso ).
Considerando -se, separadamente, as regiões brasileiras, verifica -se que, na região Norte, as universidades participaram com 17 cursos avaliados, aos quais foram atribuídos os seguintes conceitos : 1 a um curso ; 2 a três cursos ; 3 a seis cursos ; 4, o conceito modal, a sete cursos ; e 5 a nenhum curso.
Nenhum curso ficou sem conceito ( SC ).
Os centros universitários da região Norte foram representados por 25 cursos, dos quais um curso obteve conceito 1 ; 16 cursos, conceito 2, o valor modal ; seis cursos, conceito 3 ; um curso, o conceito 4 ; e nenhum curso, o conceito 5.
Um curso ficou sem conceito ( SC ).
As faculdades participaram com 21 cursos na região Norte, os quais receberam os conceitos 1 ( nenhum curso ), 2 ( 14 cursos ), 3 ( seis cursos ), 4 ( nenhum curso ) e 5 ( nenhum curso ).
Um curso ficou sem conceito ( SC ).
Os CEFET/IF participaram com cinco cursos, aos quais foram atribuídos os seguintes conceitos : 1, 2, 4 e 5 a nenhum curso ; e 3 a cinco cursos.
Nenhum curso ficou sem conceito ( SC ).
Na região Nordeste, as universidades participaram com 49 dos 182 cursos da área de Engenharia Civil oferecidos.
O conceito modal foi 5, atribuído a 15 cursos.
Outros 14 cursos obtiveram o conceito 3 ; 13 cursos, o conceito 4 ; cinco cursos, o conceito 2 ; e um curso, o conceito 1.
Um curso ficou sem conceito ( SC ).
Os centros universitários \textbf{contaram} com 64 cursos participantes na região Nordeste, 31 dos quais receberam o conceito modal 2.
Para os demais cursos, foram atribuídos os conceitos 3 ( 20 cursos ), 4 ( quatro cursos ), 1 ( dois cursos ) e 5 ( um curso ).
Seis cursos oferecidos por centros universitários no Nordeste ficaram sem conceito ( SC ).
As faculdades foram representadas por 54 cursos na região Nordeste, os quais receberam os conceitos 1 ( dois cursos ), 2 ( 36 cursos ), 3 ( 12 cursos ), 4 ( um curso ), 5 ( nenhum curso ) e três cursos ficaram sem conceito ( SC ).
Na região Nordeste, os CEFET/IF participaram com 15 cursos, aos quais foram atribuídos os seguintes conceitos : 1 e 2 a nenhum curso ; 3 a dois cursos ; 4 a 12 cursos ; e 5 a um curso.
Nenhum curso ficou sem conceito ( SC ).
Na região Sudeste, as universidades concentraram 133 dos 394 cursos de Engenharia Civil da região.
Dos cursos oferecidos em IES com este tipo de organização acadêmica, nessa região, o conceito modal foi 3, atribuído a 52 cursos.
Os demais cursos receberam os conceitos 2 ( 48 cursos ), 5 ( 17 cursos ), 4 ( 15 cursos ) e 1 ( um curso ).
Nenhum curso ficou sem conceito ( SC ).
Os centros universitários participaram com 137 cursos na região Sudeste, dos quais 63 obtiveram o conceito modal 2.
O conceito 3 foi atribuído a 51 cursos ; o conceito 4 a seis cursos ; o conceito 1 a seis cursos ; e o conceito 5 a um curso.
Dez cursos ficaram sem conceito ( SC ) nessa combinação de organização acadêmica e região.
Já as faculdades foram representadas por 110 cursos na região Sudeste, dos quais sete ficaram sem conceito ( SC ).
O conceito modal 2 foi atribuído a 64 cursos.
Aos demais cursos foram atribuídos os conceitos 3 ( 21 cursos ), 1 ( 11 cursos ), 4 ( quatro cursos ) e 5 ( três cursos ).
Os CEFET/IF participaram com 14 cursos na região Sudeste, aos quais foram atribuídos os seguintes conceitos : 1, 2 e 3 a nenhum curso ; 4 a 12 cursos ; e 5 a dois cursos.
Nenhum curso ficou sem conceito ( SC ).
Dos 180 cursos da região Sul, 77 eram mantidos por universidades, para os quais o conceito modal foi 3, atribuído a 32 cursos.
Os demais cursos receberam os conceitos 4 ( 20 cursos ), 2 ( 17 cursos ), 5 ( oito cursos ) e 1 ( nenhum curso ).
Em IES com esse tipo de organização acadêmica, nenhum curso da região Sul ficou sem conceito ( SC ).
Dos 59 cursos participantes de centros universitários da região Sul, três ficaram sem conceito ( SC ), e 29 cursos receberam conceito 3, o valor modal.
Os conceitos \textbf{restantes} ficaram assim distribuídos : 2 ( 21 cursos ), 4 ( quatro cursos ), 1 ( dois cursos ) e 5 ( nenhum curso ).
Quanto aos 40 cursos vinculados a faculdades na região Sul, 18 cursos receberam o conceito 2, valor modal.
Os demais conceitos ficaram assim distribuídos : 3 ( 13 cursos ), 1 ( dois cursos ), 4 ( um curso ) e 5 ( nenhum curso ).
Dos 40 cursos, seis ficaram sem conceito ( SC ).
Os CEFET/IF participaram com 4 cursos na região Sul, aos quais foram atribuídos os seguintes conceitos : 1, 2 e 5 a nenhum curso ; e 3 e 4 a dois cursos.
Nenhum curso ficou sem conceito ( SC ).
Na região Centro-Oeste, 24 dos 93 cursos oferecidos eram mantidos por universidades.
Para esse tipo de organização acadêmica, o conceito modal foi 3, atribuído a 11 cursos.
Em relação às universidades, nenhum curso da região ficou sem conceito ( SC ).
Os demais cursos receberam os conceitos 4 ( cinco cursos ), 5 ( quatro cursos ), 2 ( quatro cursos ) e 1 ( nenhum curso ).
Os centros universitários da região Centro-Oeste \textbf{contaram} com 29 cursos : 14 foram avaliados com o conceito 2 ( modal ) ; 11, com o conceito 3 ; dois, com o conceito 4 ; nenhum, com o conceito 5 ; e nenhum obteve o conceito 1.
Dois cursos ficaram sem conceito ( SC ).
Dos 31 cursos oferecidos em faculdades na região Centro-Oeste, 20 receberam o conceito 2 ( modal ).
Os conceitos \textbf{restantes} ficaram assim distribuídos : 3 ( seis cursos ), 4 ( dois cursos ), 1 ( dois cursos ) e 5 ( nenhum curso ).
Nessa região, um curso oferecido em faculdades ficou sem conceito ( SC ).
No que se refere aos CEFET/IF da região Centro-Oeste, nove cursos foram avaliados : seis cursos deles com o conceito 3 ( modal ) ; dois cursos com o conceito 4 ; um curso com o conceito 2 ; nenhum curso com o conceito 5 ; nenhum curso obteve o conceito 1 ; e nenhum curso ficou sem conceito ( SC ).
De maneira semelhante à forma como se \textbf{procedeu} em relação à seção anterior, esta seção encerra -se com o Gráfico 4.3, que apresenta a distribuição do Conceito Enade dos cursos de Engenharia Civil segundo a organização acadêmica da IES.
No Gráfico 4.3, apresenta -se a distribuição do Conceito Enade segundo a organização acadêmica da IES.
A leitura desse último gráfico é semelhante à dos anteriores, seguindo -se a lógica de que \textbf{poligonais} deslocadas para a esquerda correspondem aos conceitos mais baixos ( 1, 2 e 3 ) e \textbf{poligonais} deslocadas para a direita correspondem a Conceitos Enade mais altos.
Nas universidades ( linha azul ), todos os cursos obtiveram conceito 5 ou abaixo disso, já que a \textbf{poligonal} \textbf{cumulativa} alcança 100 \% para o conceito 5.
Além disso, os cursos de CEFET/IF apresentam uma \textbf{poligonal} mais à direita do que as demais, \textbf{denotando} uma melhor distribuição de Conceitos Enade.
No outro extremo, encontram -se os cursos mantidos por faculdades ( linha vermelha ), apresentando as respectivas \textbf{poligonais} mais à esquerda e, consequentemente, uma distribuição de Conceitos Enade com os menores valores.
Os cursos mantidos por universidades ( linha azul ) e centro universitários ( linha verde ) apresentam \textbf{poligonais} mais centrais e, consequentemente, uma distribuição de Conceitos Enade com os valores intermediários.

% matched lemmas: abscissa, contar, cruzar, cumulativo, denotar, equivaler, finalidade, indicado, ministrar, ofertar, ordenado, panorama, poligonal, pontilhado, proceder, reiterativo, restante, roxa, segunda, tabular, tipologia, variável, visualizar
\end{textsample}
